\phantomsection
\section*{Opening: Reading a Difficult Page with a Pencil}
\addcontentsline{toc}{section}{Opening: Reading a Difficult Page with a Pencil}

Open a page of St.\ Thomas Aquinas or Aristotle, and the prose greets you as a sea of authoritative assertion. \emph{Aristotle says. Augustine says. Act as act. Potency as potency. Being is said in many ways. Nothing is moved except by another.} Every sentence arrives in the same steady, confident register, formed by an intellectual training that few modern readers ever receive.

The real obstacle on that page is rarely unfamiliar Latin or technical vocabulary. A reader can look up \emph{potency} or \emph{accident} in two minutes and remain just as stuck as before. The real difficulty is that within a single paragraph of classical metaphysics, eight completely different kinds of sentence look identical on the page: a definition fixing how a word will be used; an unpacking of a technical vocabulary; a logical consequence of what was just granted; a substantive claim about reality; a theorem demonstrated in another treatise; a theological premise held on divine revelation; an appeal to the standing of a great author; and an actual demonstration. They share the same crisp syntax, often inside one sentence.

When you cannot tell those eight things apart, you are left with two bad options. You can accept the whole paragraph on the author's credit---which turns philosophy into deference and leaves you unable to answer a single serious question about what you accepted. Or you can dismiss the whole paragraph as medieval verbal machinery---which destroys your access to whatever genuine truth about reality the text actually uncovered. Neither option serves a Catholic reader who wants to know what is true.

\subsection*{One Competence, Six Questions}

Everything that follows trains a single practical skill: given a difficult metaphysical passage, identify exactly what is being claimed, why you are supposed to believe it, where the substantive philosophical burden lies, what actually follows if you grant it, and what happens if you say no.

To exercise that skill, you need no formal system, no special symbols, and no glossary to memorize before opening a text. You need a pencil and six ordinary English questions asked of any passage that claims to argue for something:

\begin{itemize}
\item \textbf{Claim.} What exactly is being asserted? State the claim in the author's own words with every qualification intact, and then write a \emph{plain-language reading} that drops unnecessary school diction without flattening a single distinction.
\item \textbf{Meaning.} What do the load-bearing words mean here? Fix the technical senses, the distinctions the argument turns on, and any term whose sense shifts between one occurrence and the next---and fix them \emph{before} evaluating the argument.
\item \textbf{Warrant.} Why am I supposed to accept each important premise? Go premise by premise, sorting what the text actually offers into eight ordinary kinds: \emph{definition}, \emph{logical consequence}, \emph{experience}, \emph{prior result}, \emph{metaphysical principle}, \emph{hypothesis}, \emph{revelation}, or \emph{provenance only} (the diagnosis that nothing has been offered on the page except who said it).
\item \textbf{Burden.} Which sentence is doing the real philosophical work? Locate the \emph{burden-bearing premise}: the first proposition here whose truth is not supplied merely by vocabulary or deferred to a prior result.
\item \textbf{Consequence.} What has actually been established if I grant it? State what follows in the weakest form the argument genuinely licenses, and then state conspicuously what is \emph{not yet established}---naming the specific further claims commonly made at this point that this passage has not reached.
\item \textbf{Denial.} What follows if I reject the burden-bearing premise? Classify what saying no actually costs among five possibilities: \emph{contradiction}, \emph{conflict with experience}, \emph{explanatory cost}, \emph{rival framework}, or \emph{non-assent}.
\end{itemize}

Those six questions are an order of exposition, not a rigid assembly line. In live reading you will often locate the burden first and only then realize which word under Meaning needs sharpening, or reach Consequence and discover that your plain-language reading dropped a qualifier. Iteration is normal; presenting the finished reading in order from Claim to Denial simply keeps your analysis legible to yourself and to another reader. Nor are these six headings a form to fill in or a scorecard that outputs a grade. A good analysis is continuous prose; if you deleted all six headings, the paragraphs should still read smoothly from start to finish.

\subsection*{What Governs the Reading}

One central distinction runs beneath all six questions: \emph{unpacking our concepts is not discovering something about reality}. Careful definitions are indispensable in philosophy, and conceptual clarifications are often genuinely valuable. What ruins metaphysical reading is the silent transition---the moment an argument stops explaining its words and begins telling you how the world is, without announcing the crossing and without offering an argument for it. Your automatic test on every page is therefore simple: \emph{could an intelligent opponent accept every definition offered here and still deny the substantive conclusion?} If not, the argument has unpacked a vocabulary rather than made a discovery about reality.

A second discipline keeps the water clear right beside the first: \emph{source is not warrant}. That St.\ Thomas Aquinas or Aristotle wrote a sentence proves that he held it, and gives any serious thinker good reason to read him with care and look for his argument; it does not prove that the sentence is true. The scholastic article itself honors this rule, routinely citing Aristotle, St.\ Augustine, or St.\ John of Damascus in objections that the author goes on to qualify or reject. Likewise, when reading Catholic theologians such as Aquinas, \emph{a revealed premise stays labelled}: it is never quietly dressed up as a result of natural reason, and it is never faulted for lacking a philosophical proof it never claimed.

Prose remains the canonical medium of metaphysical reasoning throughout. Symbolic notation enters only on the rare occasions where a line or two of logic exposes a specific scope ambiguity, modal shift, identity claim, or invalid inference more plainly than a paragraph can---and every displayed formula is accompanied by a sentence saying exactly what it makes easier to see. No metaphysical notion such as act, potency, essence, existence, or analogy is turned into a predicate letter merely to make a page look rigorous. Finally, when an argument rejects a rival position, that rival must appear in the strongest form its ablest defender would accept; training yourself on weak opponents leaves you helpless against real objections.

\subsection*{From Reading to Demonstration}

Most introductions to logic begin by asking the student to construct formal proofs and treat reading real philosophical texts as an afterthought. That order fails in metaphysics. A reader who cannot yet tell a definition from a metaphysical principle, or a citation from a reason, will build arguments in which both look identical and never notice. Accurate reading is load-bearing; judging or attempting a complete demonstration is what accurate reading makes possible.

We begin, therefore, by recovering what Aristotle and St.\ Thomas actually meant by \emph{demonstration}, \emph{dialectic}, and showing \emph{that} versus showing \emph{why}. From there the six questions are installed one by one on classical passages, tested against the limits of Aristotelian demonstration, and exercised on progressively harder unmarked prose---including thinkers outside the Thomistic line---before turning to full reconstructions where an attempted demonstration must be judged honestly and, when it falls short, pared back to what it genuinely proves.

One practical test governs the whole enterprise: if you have to remember a specialized apparatus of ours in order to interrogate an unmarked page of Aquinas with a pencil, the method has failed. Keep the six ordinary questions, keep your honesty about what a page proves and what it leaves open, and throw away everything else.
